\documentclass[11pt]{article}
\usepackage[T1]{fontenc}
\usepackage{lmodern,amsmath,amssymb,amsthm}
\usepackage[margin=26mm]{geometry}
\usepackage[colorlinks=true,linkcolor=blue,citecolor=blue,urlcolor=blue]{hyperref}
\newcommand{\dd}{\mathrm d}
\newcommand{\R}{\mathbb R}
\newcommand{\atanh}{\operatorname{artanh}}
\newtheorem{theorem}{Theorem}
\setlength{\emergencystretch}{2em}
\title{Proof of the local particle-number and energy laws}
\author{Companion to \emph{Particle number and energy from the shape of relativistic mirror motion}}
\date{18 November 2026}
\begin{document}
\maketitle
\section{Notation and starting identities}
Laboratory time is $t$, the mirror position is $z(t)$, and its velocity
is $v(t)=\dot z(t)$. We use $c=\hbar=1$ and count radiation in both
exterior regions for one real massless scalar field. For $t\ne t'$,
define the chord velocity and a velocity-difference functional by
\begin{equation}
 \bar v(t,t')=\frac{z(t)-z(t')}{t-t'},\qquad
 G[v]=\iint_{\R^2}\left[\frac{v(t)-v(t')}{t-t'}\right]^2\dd t\dd t'.
 \label{eq:G}
\end{equation}
The diagonal values are defined by continuity. The parameter
$\mu=(\omega-\omega')/(\omega+\omega')\in(-1,1)$ is the frequency
asymmetry, with $\omega,\omega'>0$ the outgoing and incoming frequencies.
The number and energy identities established in Appendix~A of the
accompanying manuscript~\cite{MirrorCountsManuscript} are our starting
point:
\begin{equation}
 N=\frac1{8\pi^2}\int_{-1}^1(1-\mu^2)\dd\mu
 \iint_{\R^2}
 \frac{\left[\dfrac{v(t)-v(t')}{t-t'}\right]^2\dd t\dd t'}
 {[1-\mu v(t)][1-\mu v(t')][1-\mu\bar v(t,t')]^2},
 \label{eq:pn-positive-number_appendix}
\end{equation}
\begin{equation}
 E=\frac1{6\pi}\int_\R\frac{\dot v(t)^2}{[1-v(t)^2]^3}\dd t.
 \label{eq:energy-positive-functional}
\end{equation}
Here $N$ is the total particle number in the incoming vacuum and $E$
is the corresponding total energy. The boundary conditions in the
theorem identify the second formula with the ordinary integrated
stress flux. Both displayed integrands are nonnegative.

\section{Main theorem}

\begin{theorem}[Separated fast events]
\label{thm:local}
Let $v_\lambda\in C^1(\R)$ be a family of velocities with
$|v_\lambda|<1$, $v_\lambda(t)\to0$ as $t\to\pm\infty$, and
$G[v_\lambda]\le M_G$ uniformly, where $G$ is defined in
Eq.~\eqref{eq:G} with laboratory time as its argument.
Suppose there are a fixed number $J$ of event centers
$t_j(\lambda)$ in a fixed compact set and a fixed $\rho>0$ such that
the intervals $I_j=(t_j-\rho,t_j+\rho)$ are disjoint. Outside their
union, require $|v_\lambda|\le1-c$ with fixed $c>0$.

On each full interval, suppose
\begin{equation}
 \sigma_j v_\lambda(t_j+\Delta t)
 =V_j-a_j(\Delta t)^{2m_j}-R_j(\Delta t),
 \label{eq:profile}
\end{equation}
where $|\Delta t|<\rho$, $\sigma_j\in\{-1,1\}$,
$V_j\to1$, $m_j\ge1$ is a fixed integer, and
$0<a_-\le a_j\le a_+<\infty$. For a fixed $\alpha>0$, assume
\begin{equation}
 \begin{aligned}
 |R_j(\Delta t)|&\le M|\Delta t|^{2m_j+\alpha},\\
 |R_j'(\Delta t)|&\le M|\Delta t|^{2m_j+\alpha-1},
 \end{aligned}
 \label{eq:remainders}
\end{equation}
and, throughout $I_j$,
\begin{equation}
 \begin{aligned}
 \sigma_j v_\lambda(t_j+\Delta t)&\ge\tfrac12,\\
 1-\sigma_j v_\lambda(t_j+\Delta t)
     &\ge1-V_j+\tfrac12a_j|\Delta t|^{2m_j}.
 \end{aligned}
 \label{eq:coercive}
\end{equation}
All constants are independent of $\lambda$ and $j$. Define the peak
Lorentz factors and nonnegative peak rapidities by
\begin{equation}
 \gamma_j=(1-V_j^2)^{-1/2}=\cosh\eta_j,
 \qquad \eta_j=\atanh V_j.
 \label{eq:local-heights}
\end{equation}
Then the particle number in Eq.~\eqref{eq:pn-positive-number_appendix}
satisfies
\begin{align}
 N&=\sum_{j=1}^J C_{m_j}\log\gamma_j
       +o\!\left(\sum_{j=1}^J\log\gamma_j\right)\nonumber\\
  &=\sum_{j=1}^J C_{m_j}\eta_j
       +o\!\left(\sum_{j=1}^J\eta_j\right).
 \label{eq:localN}
\end{align}
The coefficient is
\begin{equation}
 C_m=\frac1{\pi m}\sum_{p=1}^m\cot\frac{\pi p}{2m+1}.
 \label{eq:Cm}
\end{equation}
The limits $V_j\to1$ may occur at unequal rates.

For energy, assume in addition $v_\lambda\in C^2(\R)$,
the boundary limits
\begin{equation}
 \lim_{t\to\pm\infty}
 \frac{\dot v_\lambda(t)}{[1\mp v_\lambda(t)]^2[1\pm v_\lambda(t)]}=0
 \quad\text{for both sign choices},
 \label{eq:stress-boundaries}
\end{equation}
and the uniform bound
\begin{equation}
 \int_{\R\setminus\bigcup_j I_j}
       \frac{\dot v_\lambda^{\,2}}{(1-v_\lambda^2)^3}\dd t\le M_E.
 \label{eq:outer-energy-bound}
\end{equation}
Then
\begin{equation}
 \begin{aligned}
 E&\sim\sum_{j=1}^J D_{m_j}a_j^{1/(2m_j)}
                              \gamma_j^{2+1/m_j}\\
  &=\sum_{j=1}^J D_{m_j}a_j^{1/(2m_j)}
                              \cosh^{2+1/m_j}\eta_j,
 \end{aligned}
 \label{eq:localE}
\end{equation}
where
\begin{equation}
 D_m=\frac{2^{1/(2m)}[1-1/(2m)]}{24\sin[\pi/(2m)]}.
 \label{eq:local-energy-coefficient}
\end{equation}
\end{theorem}

Here $\Delta t=t-t_j$ is a displacement in laboratory time, and
$R_j$ is the remainder after subtracting the peak value and leading
even power. The lower bound~\eqref{eq:coercive} excludes an additional
fast maximum in the same interval. These are uniform assumptions on
a family, rather than a Taylor expansion of one fixed trajectory.

\section{Proof for one positive-velocity event}
\subsection{Characteristic time scale and logarithmic number growth}
\label{app:local-scaling}

Take a single event with $\sigma=+1$, translate its center to zero,
and suppress the family index. In physical laboratory time its
profile is
\begin{equation}
 v(t)=V-a t^{2m}-R(t),\qquad |t|<\rho,\qquad V\to1.
 \label{eq:single-event-profile}
\end{equation}
Only the endpoint $\mu=1$ can be enhanced by this event.
Write its distance as $\delta\mu=1-\mu$ and choose a fixed small
$\delta\mu_0>0$. For $0<\delta\mu<\delta\mu_0$, set
\begin{equation}
 \mu=1-\delta\mu,\qquad
 \ell=\left(\frac{1-\mu V}{a}\right)^{1/(2m)}.
 \label{eq:local-scale}
\end{equation}
The exact denominator at the peak is
$1-\mu V=1-V+V\delta\mu$.
The coefficient $a$ has units of time to the power $-2m$, so
$\ell$ has units of laboratory time. For $\mu$ near one, it is the
characteristic width of the region where the local speed variation
$a|t|^{2m}$ competes with the phase-rate gap $1-\mu V$.
As $V\to1$ and $\mu\to1$, the uniform lower bound on $a$ makes this
width shrink to zero.


In the time integral over $(-\rho,\rho)^2$, substitute
$t\mapsto\ell t$ and $t'\mapsto\ell t'$.
For readability the dimensionless dummy variables after this
substitution are again denoted by $t,t'$. For example,
\begin{equation}
 \frac{1-\mu v(\ell t)}{1-\mu V}
 =1+\mu t^{2m}+\frac{\mu R(\ell t)}{1-\mu V}
 \longrightarrow1+t^{2m}.
 \label{eq:scaled-denominator}
\end{equation}
The remainder divided by $1-\mu V=a\ell^{2m}$ is
$O(\ell^\alpha|t|^{2m+\alpha})$. The derivative bound controls
divided differences in the same way, including $t=t'$.
The chord velocity is the average of $v$ over the connecting time
interval. Its denominator, divided by $1-\mu V$, tends to
$1+Q_{2m}(t,t')$, where
\begin{align}
 Q_{2m}(t,t')
 &=\int_0^1[(1-\theta)t+\theta t']^{2m}\dd\theta\nonumber\\
 &=\frac{t^{2m+1}-(t')^{2m+1}}{(2m+1)(t-t')}.
 \label{eq:local-Q}
\end{align}
The second form is continued at $t=t'$.

The velocity difference contributes a factor $(1-\mu V)^2$,
while the four denominator powers contribute $(1-\mu V)^4$.
The two time measures cancel the two powers of $\ell$ in the
squared time difference. Consequently the original double-time
integral over this event square is
\begin{equation}
 \frac{K_{2m}+o(1)}{(1-\mu V)^2},
 \label{eq:inner-limit}
\end{equation}
with the positive, dimensionless coefficient
\begin{equation}
 \begin{aligned}
 K_{2m}={}&\iint_{\R^2}
 \left[\frac{t^{2m}-(t')^{2m}}{t-t'}\right]^2\\
 &\times\frac{\dd t\dd t'}
 {(1+t^{2m})[1+(t')^{2m}][1+Q_{2m}(t,t')]^2}.
 \end{aligned}
 \label{eq:K}
\end{equation}
In particular, the coefficient $a$, which sets the physical time
scale, cancels from the particle count.

We justify the limit before doing the $\mu$ integral.
The nonnegative homogeneous polynomial $Q_{2m}$ vanishes only when
$t=t'=0$: otherwise the line segment in its defining integral
cannot have zero even power throughout. Compactness on
$|t|^{2m}+|t'|^{2m}=1$ therefore gives
$Q_{2m}\asymp |t|^{2m}+|t'|^{2m}$.
For $\mu$ sufficiently close to one, the coercive bound in the
theorem controls the scaled point and chord denominators by these
same positive polynomials. The derivative remainder bounds the
divided difference by $C(|t|+|t'|)^{2m-1}$.
Thus a constant times
\begin{equation}
 \frac{(|t|+|t'|)^{4m-2}}
 {(1+|t|^{2m})(1+|t'|^{2m})
                  (1+|t|^{2m}+|t'|^{2m})^2}
 \label{eq:majorant}
\end{equation}
majorizes the scaled integrand, extended by zero outside its
expanding event square. This function is bounded near the origin.
On $|t|\ge\max(1,|t'|)$ it is at most
$C|t|^{-2m-2}/(1+|t'|^{2m})$, an integrable bound; exchange the
variables for the other region. Dominated convergence proves
Eq.~\eqref{eq:inner-limit} jointly as $V\to1$ and $\mu\to1$,
with no restriction on their relative rates.

The weight is $1-\mu^2=2\delta\mu-(\delta\mu)^2$.
The leading event contribution to $N$ is therefore
\begin{equation}
 \frac{K_{2m}}{4\pi^2}
 \int_0^{\delta\mu_0}
       \frac{\delta\mu\,\dd(\delta\mu)}{(1-V+V\delta\mu)^2}.
 \label{eq:log-integral}
\end{equation}
This last integral is elementary:
\begin{align}
 &\int_0^{\delta\mu_0}
       \frac{\delta\mu\,\dd(\delta\mu)}{(1-V+V\delta\mu)^2}
       \nonumber\\
 &\quad=\frac1{V^2}\left[
       \log\frac{1-V+V\delta\mu_0}{1-V}
       +\frac{1-V}{1-V+V\delta\mu_0}-1\right]\nonumber\\
 &\quad=\log\frac1{1-V}+O(1).
 \label{eq:log}
\end{align}
The omitted $(\delta\mu)^2$ weight gives a bounded contribution.
To control the $o(1)$ in Eq.~\eqref{eq:inner-limit}, fix a small
threshold for $1-\mu V$. Below it, uniform convergence bounds the
error by an arbitrarily small multiple of the integral in
Eq.~\eqref{eq:log}; above it, the denominators have a fixed gap and
$G[v]$ bounds the contribution. Sending $V\to1$ first and then the
threshold to zero gives $o(\log[1/(1-V)])$.
The $\mu$ range away from $\mu=1$ is bounded for the same reason.

To complete the single-event localization, we must also control time
pairs outside $(-\rho,\rho)^2$. Here and below $t,t'$ again denote
physical laboratory times. Outside $(-\rho/2,\rho/2)$, the
coercive bound and the outer speed gap give $1-|v|\ge c_1>0$.
For a pair outside the full event square, its connecting interval
also has a uniform average gap. If its length is at least $4\rho$,
at least half lies outside the full event interval. If its length
is shorter and it enters the smaller core, it must traverse an
annulus of width $\rho/2$ before reaching the endpoint outside the
full interval. If it misses the core, the gap $c_1$ applies
throughout. In all cases $1-|\bar v(t,t')|\ge c_0>0$.

When both endpoints avoid the core, all denominator factors are
uniformly bounded below, and $G[v]\le M_G$ bounds their contribution.
When one endpoint is in the core, its partner is at least
$\rho/2$ away. The elementary inequality
\begin{equation}
 \frac{(1-\mu^2)(v-w)^2}{(1-\mu v)(1-\mu w)}\le4,
 \qquad |v|,|w|<1,
 \label{eq:cross}
\end{equation}
then bounds the full kernel by $C/(t-t')^2$.
To verify the inequality, take $v\ge w$; the ratio increases with
$v$ and decreases with $w$, as the derivative numerators contain
$2-\mu(v+w)>0$. Its supremum is its value $4$ at $v=1,w=-1$.
The partner-time integral over $|t-t'|\ge\rho/2$ is finite and the
core has bounded length. All these remaining pairs contribute
$O(1)$, so they cannot change the leading logarithm.

\subsection{The coefficient of the single-event number law}
Equations~\eqref{eq:log-integral} and \eqref{eq:log} already show that
the coefficient of $\log\gamma$ is $K_{2m}/(2\pi^2)$, independently
of the time scale $a$. To evaluate it, take the exact flight
\begin{equation}
 t(z)=-\frac zV-\frac{z^{2m+1}}{(2m+1)V^{2m+1}\tau^{2m}},
 \qquad 0<V<1,\quad \tau>0.
 \label{eq:family}
\end{equation}
It has one negative-velocity fast event. Reflection $z\mapsto-z$,
$\mu\mapsto-\mu$ leaves the number kernel invariant, so the
positive-event argument just proved applies. With
$\xi=-z/(V\tau)$, one has
$t=\tau[\xi+\xi^{2m+1}/(2m+1)]$ and
$v=-V/(1+\xi^{2m})$. Inversion near zero gives
\begin{equation}
 -v(t)=V-\frac V{\tau^{2m}}t^{2m}+O(t^{4m}),
 \qquad a=\frac V{\tau^{2m}}.
 \label{eq:flight-local-profile}
\end{equation}
The differentiated remainder has order $O(|t|^{4m-1})$ uniformly
for $V$ near one. A sufficiently small fixed event interval supplies
the coercive bound; outside it the fixed shape has a speed gap.
The velocity is $V$ times a fixed smooth function of $t/\tau$ with
tail $O(|t|^{-2m/(2m+1)})$, whose derivative has one more inverse
power. The fixed shape and its derivative are square integrable.
Parseval and Cauchy--Schwarz therefore give finite $G$; multiplying
by $V$ preserves its uniform bound. Thus all number hypotheses
hold for this comparison family.

The exact finite-speed count established in Sec.~III and Appendix~B
of the accompanying manuscript~\cite{MirrorCountsManuscript} gives
\begin{equation}
 N_m=C_mF(\eta_*),\qquad
 F(\eta_*)=\frac{\eta_*}{\tanh\eta_*}-1
          =\frac{\atanh V}{V}-1,\qquad \eta_*=\atanh V.
 \label{eq:F}
\end{equation}
Since $F(\eta_*)=\log\gamma+O(1)$, comparison with the already
proved local logarithm fixes $K_{2m}/(2\pi^2)=C_m$.
The manuscript obtains that finite-speed count from its generalized-Airy
calculation, independently of the local asymptotic law. Finally,
\begin{equation}
 \log\frac1{1-V}=2\log\gamma+\log(1+V),\qquad
 \log\gamma=\eta-\log2+o(1),\quad \eta=\atanh V,
 \label{eq:local-log-gamma}
\end{equation}
establishes Eq.~\eqref{eq:localN} for one event in both forms.

\subsection{Leading energy and its characteristic time scale}

The same local profile gives a power law for energy. We now use
physical laboratory time in Eq.~\eqref{eq:single-event-profile},
and the energy functional~\eqref{eq:energy-positive-functional}.
This time choose
$\ell=[(1-V)/a]^{1/(2m)}$ and substitute $t\mapsto\ell t$.
This is the characteristic laboratory-time width over which the
speed falls below its peak by an amount comparable to $1-V$.
It is the $\mu=1$ member of the scales in Eq.~\eqref{eq:local-scale}.
Again denote the dimensionless integration variable by $t$.
On compact sets the profile and its differentiated remainder give
\begin{align}
 1-v(\ell t)^2
    &=2(1-V)(1+t^{2m})[1+o(1)],\nonumber\\
 \dot v(\ell t)
    &=-2m a\ell^{2m-1}t^{2m-1}
                          +o(a\ell^{2m-1}).
 \label{eq:energy-scaled-profile}
\end{align}
The first relation is uniform on compact sets; the derivative
remainder vanishes there by Eq.~\eqref{eq:remainders}.
On the full expanding event interval, the coercive and derivative
bounds majorize the scaled energy integrand by
$C|t|^{4m-2}/(1+|t|^{2m})^3$.
This is integrable at zero and decays as $|t|^{-2m-2}$ at infinity.
Dominated convergence therefore gives the event contribution
\begin{align}
 E_{\mathrm{event}}\sim{}&\frac{(2m)^2}{48\pi}
        a^{1/(2m)}(1-V)^{-1-1/(2m)}\nonumber\\
 &\times\int_\R\frac{t^{4m-2}}{(1+t^{2m})^3}\dd t.
 \label{eq:betaenergy}
\end{align}
Here $E_{\mathrm{event}}$ denotes the energy integral restricted to
the single event interval.

Set $h=t^{2m}$ on the positive half-line. The beta integral
\cite[Eqs.~5.12.1 and 5.12.3]{DLMF}, followed by the gamma recurrence and reflection
identities~\cite[Eqs.~5.5.1 and 5.5.3]{DLMF}, gives
\begin{align}
 \int_\R\frac{t^{4m-2}}{(1+t^{2m})^3}\dd t
 &=\frac1m B\left(2-\frac1{2m},1+\frac1{2m}\right)\nonumber\\
 &=\frac{\Gamma(2-1/(2m))\Gamma(1+1/(2m))}{2m}
       \nonumber\\
 &=\frac{\pi[1-1/(2m)]}{4m^2\sin[\pi/(2m)]}.
 \label{eq:local-energy-beta}
\end{align}
In the last step,
$\Gamma(2-1/(2m))=[1-1/(2m)]\Gamma(1-1/(2m))$,
$\Gamma(1+1/(2m))=\Gamma(1/(2m))/(2m)$, and
$\Gamma(u)\Gamma(1-u)=\pi/\sin(\pi u)$.
Substitution into Eq.~\eqref{eq:betaenergy} yields
\begin{align}
 E_{\mathrm{event}}&\sim\frac{1-1/(2m)}{48\sin[\pi/(2m)]}
          a^{1/(2m)}(1-V)^{-1-1/(2m)}\nonumber\\
    &\sim D_m a^{1/(2m)}\gamma^{2+1/m}
     =D_m a^{1/(2m)}\cosh^{2+1/m}\eta.
 \label{eq:one-event-energy}
\end{align}
We used $\gamma^{-2}=(1-V)(1+V)$ and $V\to1$.
The uniformly bounded outer energy in
Eq.~\eqref{eq:outer-energy-bound} is subleading.

For the exact-flight normalization in Eq.~\eqref{eq:flight-local-profile},
$a=V/\tau^{2m}$, so the corresponding leading energy is
\begin{equation}
 E\sim\frac{1-1/(2m)}{48\tau\sin[\pi/(2m)]}
                 (1-V)^{-1-1/(2m)}.
 \label{eq:main-energy-coefficient}
\end{equation}
The additional factor $V^{1/(2m)}$ tends to one.
The time scale cancels from the number coefficient but remains in
the energy coefficient. This is why the number grows logarithmically
with the peak Lorentz factor, whereas energy also resolves the
width and order of the peak.

\section{Negative velocities and several separated events}

For a single event with $\sigma=-1$, reflect the worldline:
$z\mapsto-z$, $v\mapsto-v$, and change $\mu\mapsto-\mu$.
The number kernel, its $\mu$ measure, and the energy functional
are invariant. The enhanced endpoint is now $\mu=-1$, with
$\delta\mu=1+\mu$. The preceding proof applies unchanged to the
reflected positive event. Both coefficients therefore depend on
the local peak order and magnitude, not on its sign.

For several events, the analysis on each same-event square
$I_j\times I_j$ is exactly the single-event analysis, with
$\Delta t=t-t_j$ before rescaling. It remains to bound pairs that
do not belong to any such square, including pairs from different
fast events. We give this step explicitly because additivity is
not an assumption about loss of phase coherence.

Let $I_j^0=(t_j-\rho/2,t_j+\rho/2)$.
Outside these smaller cores, Eq.~\eqref{eq:coercive} and the outer
speed gap give a uniform $1-|v|\ge c_1>0$.
Consider a connecting interval for a pair outside all the
same-event squares. If its length is at least $4J\rho$, at least
half of it lies outside the union of the full event intervals,
whose total length is $2J\rho$; its average gap is at least $c/2$.
For a shorter interval that misses every core, use $c_1$.
If it meets a core, its endpoints cannot both lie in that core's
full interval, because the pair is outside every same-event
square. It therefore crosses an annulus of width at least
$\rho/2$. Dividing its positive gap integral on this annulus by
the upper length bound $4J\rho$ gives another fixed positive
average gap. Thus
\begin{equation}
 1-|\bar v(t,t')|\ge c_0>0
 \label{eq:outer-chord-gap}
\end{equation}
uniformly for all these pairs.

If both endpoints avoid the smaller cores, the point denominators
are also bounded below, so their total contribution is at most
$C G[v_\lambda]$. Otherwise one endpoint belongs to a bounded
core and its partner is at least $\rho/2$ away. Apply
Eq.~\eqref{eq:cross} and the chord gap to obtain the bound
$C/(t-t')^2$. Integrating over the separated partner and then over
the finite union of cores gives $O(1)$, uniformly in the family.
This establishes boundedness of all outer and cross-event number
contributions without estimating an oscillatory interference term.

The same-event squares give
$\sum_j[C_{m_j}\log\gamma_j+o(\log\gamma_j)]$.
For a fixed finite number of events, the maximum of the finitely
many relative errors tends to zero. Their sum is therefore
$o(\sum_j\log\gamma_j)$, even when the $V_j$ approach unity at
different rates. Equation~\eqref{eq:local-log-gamma} supplies the
rapidity statement. For energy the exact functional is already a
single integral, so the event intervals add directly.
Its complement is bounded by Eq.~\eqref{eq:outer-energy-bound},
and the same finite-sum error argument proves
Eq.~\eqref{eq:localE}. This completes the theorem.

\section{Scope of the result}

The proof gives the stated little-$o$ number remainder; a bounded
remainder is not established for the entire class. Peaks that
remain uniformly subluminal can be included in the bounded outer
part. Every additional peak that tends to light speed must be
retained, even if its Lorentz factor grows more slowly than the
largest one. Merging centers, a growing number or order of peaks,
or coefficients $a_j$ tending to zero or infinity require a
different uniform analysis. A near-light-speed interval of fixed
duration violates the coercive local profile. Motion outside the
listed events can contribute at leading order if either global
bound fails. These are the reasons that smoothness and asymptotic
rest alone do not imply the local laws.

\bibliographystyle{unsrt}
\bibliography{references}
\end{document}
